% How Universal Are Concept Directions? A Controlled Stress-Test of
% Representation Engineering Under RLHF
% BlackboxNLP 2026 @ EMNLP 2026
%
% (Retitled 2026-07-07: the original working title "The Myth of Universal
% Concept Directions Under RLHF" asserted fragmentation + RLHF amplification,
% which the Qwen results did not support --- refusal is largely universal under
% the prompt-based method, and RLHF slightly *consolidates* directions. The
% question form matches the actual findings and the Repro-track framing.)
%
% Compile with: pdflatex main.tex && bibtex main && pdflatex main.tex && pdflatex main.tex

\documentclass[11pt]{article}

% === Official ACL style (acl.sty vendored in this directory) ===
% [review] adds line numbers + anonymization for submission;
% switch to \usepackage[final]{acl} for camera-ready.
\usepackage[review]{acl}

% === Standard packages ===
% NOTE: acl.sty already loads natbib and hyperref and sets the two-column
% *ACL page geometry -- do not load geometry/natbib/hyperref again.
\usepackage{times}
\usepackage{latexsym}
\usepackage[T1]{fontenc}
\usepackage[utf8]{inputenc}
\usepackage{amsmath}
\usepackage{amssymb}
\usepackage{graphicx}
\usepackage{booktabs}
\usepackage{xcolor}
\usepackage{subcaption}
\usepackage{multirow}

% === Custom macros ===
\newcommand{\cosim}{\text{cos-sim}}
\newcommand{\dmodel}{d_{\text{model}}}
\newcommand{\vdir}{\mathbf{d}}
\newcommand{\vact}{\mathbf{a}}
\newcommand{\R}{\mathbb{R}}

% === Metadata ===
\title{How Universal Are Concept Directions? \\ A Controlled Stress-Test of Representation Engineering Under RLHF}

% ACL author blocks are separated with \And / \AND (not \and); under
% [review] the block is hidden automatically for anonymity.
\author{
  Author A \\
  Affiliation \\
  \texttt{email@example.com}
  \And
  Author B \\
  Affiliation \\
  \texttt{email@example.com}
}

\begin{document}
\maketitle

% ============================================================
\begin{abstract}
Representation engineering extracts a single linear direction for a concept
such as refusal or honesty and treats it as valid wherever the concept
occurs. We stress-test this universality assumption under controlled
conditions: for each concept we compare directions extracted within four
semantic domains against a direction pooled across domains, on both the
aligned and pre-alignment checkpoints of Qwen~2.5~3B, with confounds from
mixed data sources, templated responses, and domain imbalance explicitly
removed. Universality holds for refusal: per-domain and global directions
agree closely (cosine $\geq 0.83$ at every layer) under two independent
operationalizations; a global steering direction shifts held-out
behavior as effectively as domain-specific ones; and probes transfer
across domains almost perfectly. It fails partially for honesty, where
per-domain directions diverge in early and middle layers, reconverge only
partially with depth, and resist cross-domain probing, indicating that the
divergence is not merely a difference in surface statistics. We report these
cosines against two references that a raw comparison to a pooled direction
lacks: a leave-one-domain-out global direction, since a domain contributes to
the pooled one and $k$ orthogonal domains would still score $1/\sqrt{k}$, and
a within-domain split-half noise floor estimated at matched sample size, which
separates genuine divergence from estimation error. Contrary to our pre-registered hypothesis, alignment does not
amplify this fragmentation; if anything it slightly consolidates directions.
Finally, strong refusal fragmentation that we observed in a smaller pilot
model did not replicate on the primary model under identical data,
cautioning against generalizing single-model geometry findings. Universality
is thus a property of particular concepts at particular depths in particular
models, not of concept directions as such.
\end{abstract}

% ============================================================
\section{Introduction}
\label{sec:intro}

A growing body of interpretability work holds that large language models
encode high-level concepts such as honesty, refusal, and harmfulness as
linear directions in activation space \citep{zou2023representation,
marks2023geometry}. The practical appeal is clear: if a concept is one
vector, it can be monitored with a dot product and controlled by adding that
vector back into the residual stream \citep{turner2023activation,
li2024inference}. Safety applications built on this premise, from refusal
probes to steering-based guardrails, inherit a strong implicit assumption:
that the direction extracted from one distribution of examples is the
\emph{same} direction wherever the concept occurs. A refusal direction fit
on requests about violence is presumed to govern refusals of privacy
violations; an honesty direction fit on trivia is presumed to govern honesty
about arithmetic.

This universality assumption is rarely tested, and testing it naively is
misleading. If the examples for different semantic domains come from
different datasets, differ in template structure, or differ in quantity,
then per-domain directions will diverge for reasons that have nothing to do
with the concept. Our own pilot experiment furnishes a cautionary example:
an apparent refusal fragmentation that vanished entirely on the primary
model under matched data.

We therefore stress-test universality under controlled conditions. For two
concepts, \textit{refusal} and \textit{honesty}, each decomposed into four
semantic domains, we compare per-domain concept directions against a pooled
global direction at every layer of Qwen~2.5~3B. The design removes the
obvious confounds: honesty domains are carved from a single source
(TruthfulQA), refusal is operationalized two independent ways (the standard
harmful-vs-harmless prompt contrast and a response-based contrast), domains
are balanced by capping, and confidence intervals come from a pair-level
bootstrap that re-estimates every direction per resample. Geometry is
checked against behavior with activation steering on held-out prompts, and
the whole analysis is repeated on the pre-alignment Base checkpoint to ask
whether RLHF amplifies or suppresses whatever structure exists. Two of our
three pre-registered hypotheses predicted failure of universality; we report
what we found rather than what we predicted.

We report four findings. \textbf{(1)}~Refusal directions are largely
universal: per-domain and global directions agree with cosine similarity
0.83--0.94 across all layers, under both operationalizations, and probes
trained on one domain transfer to the others at up to 0.98 accuracy
(\S\ref{sec:results-refusal}). \textbf{(2)}~Honesty is not. Early- and
middle-layer directions diverge, reconverging only partially
by the output layers, and the divergence is not merely lexical: honesty
probes fail to transfer across domains (\S\ref{sec:results-honesty}). We
qualify this finding in two ways that the raw geometry cannot. Against a
leave-one-domain-out global direction the effect is not an artifact of a
domain contributing to the pooled direction it is compared with; against a
within-domain split-half noise floor at matched sample size, part of the
early-layer divergence reflects estimation error rather than geometry
(\S\ref{sec:results-controls}). \textbf{(3)}~The
geometry is behaviorally real: a single global refusal direction steers
held-out prompts as effectively as each domain's own direction, with no
systematic per-domain advantage (\S\ref{sec:results-steering}). \textbf{(4)}~Alignment does not
fragment concept representations; the Instruct checkpoint is in fact
marginally \emph{more} universal than its Base (\S\ref{sec:results-rlhf}).
In addition, fragmentation observed on a smaller pilot model failed to
replicate under identical data (\S\ref{sec:results-pilot}).

These results neither vindicate nor refute representation engineering; they
bound it. Universality is real where it has been most prominently claimed
(refusal), fails elsewhere (honesty, early depth), and varies across models
in ways that single-model studies cannot detect. We release the full
pipeline, data manifests, and validation tooling.


% ============================================================
\section{Related Work}
\label{sec:related}

\paragraph{Linear concept representations.}
\citet{zou2023representation} introduced representation engineering as a
top-down alternative to circuit analysis, extracting concept directions by
contrasting activations on stimuli that do and do not express a concept, and
demonstrated reading and control across a broad range of concepts.
\citet{marks2023geometry} showed that true and false statements separate
linearly in middle-layer activations, and \citet{hernandez2023linearity}
found that many relational mappings are well approximated by linear
transformations. This literature generally extracts one direction per
concept per model and evaluates it in-distribution.

\paragraph{Generalization and reliability of extracted directions.}
The assumption that an extracted direction is valid beyond its extraction
distribution has been questioned before, and we do not claim otherwise.
\citet{tan2024analysing} measure the generalisation and reliability of
steering vectors across many concepts and models, and find both to be
highly variable: steerability differs sharply by concept, and vectors that
work in-distribution frequently fail out of it. On the honesty side,
\citet{levinstein2024still} show that probes trained to detect truthfulness
fail to generalize across distributions, and argue that this reflects a
conceptual difficulty in isolating truth rather than a fixable engineering
problem. Our honesty result is consistent with both findings and should be
read as converging evidence rather than as a first observation.

What our design adds is the controlled attribution. Prior work varies the
concept, the dataset, and often the model together, so a generalization
failure cannot be assigned to any one of them. We hold the concept, the
source dataset, the pair schema, and the model fixed, and vary only the
semantic domain, which turns generalization into a per-concept, per-layer,
per-model measurement with a stated within-domain noise floor
(Section~\ref{sec:controls}). We also separate the two explanations that a
raw cosine cannot distinguish: directions that genuinely differ, and
directions too noisily estimated to compare.

\paragraph{Low-dimensional structure in truth representations.}
\citet{burger2024truth} report that truth is carried by a shared
two-dimensional subspace, separating a general truth direction from a
polarity-sensitive one, and that accounting for this structure makes lie
detection robust across datasets. This offers a specific alternative reading
of divergent per-domain honesty directions: domains may share a subspace
while differing in where they sit inside it, in which case a rank-1 cosine
would report shared structure as fragmentation. We test this directly on our
own activations in Section~\ref{sec:subspace} rather than leaving it open.

\paragraph{Activation steering.}
Adding a scaled direction to the residual stream at inference time shifts
model behavior without weight updates \citep{turner2023activation}, and
targeted interventions on a small set of directions can improve truthfulness
\citep{li2024inference}. Most directly related to us,
\citet{arditi2024refusal} showed that refusal in open chat models is
mediated by a single direction whose ablation disables refusal across many
harmful instructions. Our refusal results independently support their
single-direction account with a domain-stratified design, and extend it with
the observation that the account does not automatically carry over to other
concepts or, apparently, to other models.

\paragraph{Alignment and its footprint on representations.}
\citet{wei2024assessing} demonstrate that safety alignment is brittle, in
the sense that small parameter changes suffice to remove it, which motivates
our question of what alignment actually does to concept geometry. We find its footprint on
direction dispersion to be small and, if anything, consolidating. Finally,
our non-replication across model scales echoes the concern of
\citet{lieberum2023does} that interpretability findings established on one
model need not transfer to another; we provide a concrete instance at the
level of concept-direction geometry.


% ============================================================
\section{Methods}
\label{sec:methods}

Our design asks a simple question of the representation-engineering
framework: if a concept direction is extracted from data spanning several
semantic domains, how well does it agree with directions extracted from each
domain alone? We operationalize this for two concepts, \textit{honesty} and
\textit{refusal}, each decomposed into four domains, and measure agreement
both geometrically (angular dispersion) and behaviorally (activation
steering on held-out prompts). The design controls for the confounds that
would otherwise masquerade as fragmentation: source datasets that differ
across domains, templated responses, and unbalanced domain sizes.

\subsection{Contrastive Data}
\label{sec:prompts}

\paragraph{Honesty (primary).}
Domain must be the only factor that varies between the groups we compare, so
our primary experiment draws all of its data from a single source. We
partition the 817 questions of TruthfulQA \citep{lin2022truthfulqa} into
four domains by category (\textit{factual trivia}, \textit{math},
\textit{politics/opinion}, and \textit{personal advice}), supplemented
with 30 hand-written arithmetic pairs to populate the math domain. Each item
pairs a question with its best truthful answer (positive) and a plausible
falsehood (negative); activations are read from the concatenation of the
shared question and each response, so the question itself cancels in the
contrast and the measured difference isolates the response.

\paragraph{What the honesty contrast does and does not isolate.}
This construction contrasts a true against a false answer, and that is a
weaker operationalization of honesty than the term suggests. The computation
that separates the two answers plausibly differs by domain: in
\textit{math} the contrast may largely track correct versus incorrect
arithmetic; in \textit{factual trivia} it may track successful knowledge
retrieval; in \textit{politics/opinion} and \textit{personal advice},
where there is often no single checkable fact, it may track hedging,
calibration, or stance. A truthful answer also differs from a false one in
truth value, in factual correctness, and in whatever domain-specific task the
question poses, and our contrast does not separate these.

The consequence for our claim is specific and we state it plainly: a low
cosine between two per-domain honesty directions does \emph{not} by itself
show that one concept of honesty is encoded differently across domains. It is
equally consistent with the directions picking out different underlying
computations, in which case they were never estimates of a common quantity
and their disagreement is expected rather than informative. Distinguishing
these two readings would require a contrast that holds the task constant
while varying only truthfulness, for instance a model instructed to answer
the same question honestly or deceptively, so that both sides require the
same underlying computation. We do not run that experiment here, and we
therefore frame the honesty result as fragmentation of \emph{the extracted
direction} rather than of an established single concept. Refusal is less
exposed to this objection: its four domains pose the same task, declining a
harmful request, and differ in subject matter rather than in the computation
the response requires.

\paragraph{Refusal.}
We construct the refusal contrast in two independent ways. The
\textbf{prompt-based} condition follows the now-standard recipe of
\citet{arditi2024refusal}: harmful instructions (drawn from AdvBench
\citep{zou2023universal} and MedSafetyBench \citep{han2024medsafetybench},
plus hand-written privacy requests to balance the smallest domain) are
contrasted with harmless instructions, read at the final prompt token, with
no model responses involved. This is our headline refusal condition, as it
is immune to response templating. The \textbf{response-based} condition
pairs the same harmful prompts with compliant versus refusing responses; the
refusals were programmatically diversified so that no phrasing dominates the
negative class. We treat this second condition purely as a robustness check,
since synthetic responses risk encoding surface templates rather than the
target behavior. All pairs carry a framing tag distinguishing
harmful-request refusal from over-refusal of benign requests, and the two
framings are never pooled.

\paragraph{Balance.}
Because a global direction computed over unbalanced domains is dominated by
the largest one, we cap every domain at 120 pairs at extraction time
(prompt-based refusal is balanced at roughly 91--100 prompts per domain by
construction). The full dataset comprises 1{,}720 validated pairs; an
automated validator checks schema conformity, duplicate rates, and response
diversity for every file.

\subsection{Activation Extraction}
\label{sec:extraction}

Our primary models are Qwen~2.5~3B Instruct and its pre-alignment Base
checkpoint \citep{qwen2024qwen25}; Gemma~2~2B \citep{gemmateam2024gemma2}
serves as a pilot. For each prompt we record the post-block residual stream
at every layer (36 for Qwen, 26 for Gemma) at the final token position, in
half precision with NF4 4-bit weight quantization
\citep{dettmers2023qlora} to fit commodity GPUs; we return to the possible
effects of quantization in the Limitations. Extraction metadata (source
files, schemas, pair counts) is written alongside the cached activations so
that every downstream number is traceable to its data.

\subsection{Direction Computation}
\label{sec:directions}

Following \citet{zou2023representation}, the concept direction at layer
$\ell$ is the difference of means
\begin{equation}
  \vdir_\ell \;=\;
  \frac{1}{|P|}\sum_{x \in P} \vact_\ell(x)
  \;-\;
  \frac{1}{|N|}\sum_{x \in N} \vact_\ell(x),
\end{equation}
normalized to unit length, where $P$ and $N$ are the positive and negative
activation sets and $\vact_\ell(x) \in \R^{\dmodel}$ is the residual-stream
activation. Per-domain directions $\vdir_\ell^{(d)}$ use only pairs from
domain $d$; the global direction $\vdir_\ell^{(g)}$ pools all domains after
balancing.

\subsection{Angular Dispersion and Statistics}
\label{sec:angular}

Fragmentation at layer $\ell$ is summarized by the mean cosine similarity
between each per-domain direction and the global direction,
$\overline{\cosim}_\ell = \tfrac{1}{4}\sum_d
\cos\!\big(\vdir_\ell^{(d)}, \vdir_\ell^{(g)}\big)$,
together with the corresponding mean angle. A value near~1 indicates a
single shared axis; lower values indicate domain-specific structure. We also
report the full pairwise domain-by-domain cosine matrix at the most
fragmented layer, excluding the two embedding-adjacent layers, whose
directions reflect token statistics rather than concept representations.
Confidence intervals are obtained by a \emph{pair-level} bootstrap: in each
of 1{,}000 resamples we redraw prompt pairs with replacement within every
domain, recompute all per-domain and global directions from scratch, and
recompute the dispersion statistic. This propagates sampling variability in
the underlying data, rather than resampling the four domain cosines
themselves. Because early layers of a causal language model are dominated by
lexical features, we additionally report the mean cosine over the final
third of the network, where representations feed the unembedding and
behavior is read out.

\subsection{Controls for the Dispersion Statistic}
\label{sec:controls}

A cosine between a per-domain direction and a global direction is not by
itself interpretable, for three reasons that we control explicitly.

\paragraph{The global direction contains the domain it is compared against.}
Because the global direction pools all $k$ domains, each per-domain direction
is correlated with it partly by construction. The scale of this artifact is
easy to state: $k$ equal-magnitude, mutually \emph{orthogonal} domain
directions each have cosine $1/\sqrt{k}$ with their normalized sum, which is
$0.5$ at $k{=}4$. A raw pooled-global cosine near $0.5$ is therefore
indistinguishable from complete fragmentation. We accordingly report
\textbf{leave-one-domain-out} (LODO) cosines throughout: each domain's
direction is compared against a global direction estimated from the other
three domains only, so no shared data inflates the value. Under LODO, four
orthogonal domains score $0$ rather than $0.5$. We report pooled-global
cosines only alongside LODO ones, for comparability with prior work.

\paragraph{Cross-domain cosine needs a within-domain reference.}
How similar two directions from different domains \emph{should} be depends on
how similar two independent estimates of the \emph{same} domain are. We
therefore estimate a within-domain noise ceiling by repeatedly splitting each
domain's pairs into two disjoint halves ($200$ splits), computing a direction
from each half, and recording their cosine. We then recompute cross-domain
cosines with each direction estimated from the \emph{same number of pairs} as
one half. Matching the per-estimate sample size matters because cosine
estimation error grows quickly as the number of pairs falls relative to
$d_{\text{model}}$; without the match, a small domain would look fragmented
purely for being small. We report the gap between the within- and
cross-domain distributions, which is zero when directions from different
domains agree as well as two estimates of one domain.

\paragraph{An early-layer direction may be noise.}
If a concept is not linearly encoded at a given depth, a low cosine reflects
the absence of a signal rather than the fragmentation of one. Before
interpreting any layer's geometry we therefore establish that a direction
exists there, using two tests. First, a sign-flip permutation null: we flip
the positive/negative assignment of each pair independently with probability
$0.5$ and recompute $\lVert\Delta\mu\rVert$, which preserves the pairing and
the activation marginals while destroying the systematic contrast, giving a
null distribution of $500$ draws against which the observed separation is
tested. Second, a \emph{cross-validated} projection AUC: the direction is fit
on training folds and scored only on a held-out fold ($5$-fold). The
cross-validation is essential rather than cosmetic. With
$d_{\text{model}} \geq n$, a difference-of-means direction fits the noise, and
an in-sample AUC runs far above chance even for activations with no concept
signal at all; in a synthetic check at $n{=}120$, $d{=}128$ it reaches $0.86$
on pure noise, against $0.5$ under cross-validation. Layers that fail these
tests are reported as uninterpretable rather than as fragmented.

\subsection{Shared-Subspace Analysis}
\label{sec:subspace}

Low pairwise cosines are consistent with domains sharing a low-dimensional
subspace while occupying different positions within it, the structure
\citet{burger2024truth} report for truth. We test this at three levels of
stringency. We first take the singular spectrum of the stacked per-domain
directions, summarizing it by the fraction of spectral mass at rank~$1$ and
rank~$2$ and by the participation ratio, an effective dimensionality that
requires no cutoff and runs from $1$ (one shared direction) to $k$ (mutually
orthogonal). We then compare the top-$2$ subspaces of each domain's
mean-centered per-pair contrast vectors via principal angles, which asks
whether domains excite overlapping subspaces even where their mean directions
diverge. Finally, and most stringently, we run a leave-one-domain-out
subspace test: we build a rank-$k$ subspace from the other domains' directions
and measure the squared projection of the held-out direction onto it. Since a
higher-rank subspace captures any vector more easily, we report the chance
baseline for a random rank-$k$ subspace in $d$ dimensions, $k/d$, and judge
the result against that rather than against zero.

\subsection{Cross-Domain Probing}
\label{sec:probe-method}

Diverging difference-of-means directions need not imply that the domains use
different information; the divergence could be a difference in surface
statistics that a supervised classifier would look past. To test this we
train an $\ell_2$-regularized logistic probe on one domain's activations at a
given layer (positive versus negative), standardizing features on the
training domain, and evaluate its accuracy on every other domain's held-out
pairs. Within each domain we hold out 25\% of the pairs as the test split;
the probe is trained on the remaining 75\% with $\ell_2$ regularization
($C{=}1$) and the same fixed seed used throughout, and positive and negative
activations for a pair are never split across train and test. High
off-diagonal accuracy despite low direction cosines would mean the domains
share a decodable concept subspace and the geometric fragmentation is
superficial; off-diagonal accuracy near chance, with high on-diagonal
accuracy, would mean the divergence is not merely a difference in surface
statistics.

We are deliberate about how much this test can establish. A probe trained on
one domain and evaluated on another asks whether \emph{that particular
linear decision boundary} transfers out of distribution. Poor transfer is
consistent with at least two situations: the domains may encode the concept
along genuinely different axes, or they may share a weaker common feature
that the probe declines to use because a stronger domain-specific feature
predicts the training domain better. The second is a real possibility, and
low off-diagonal accuracy therefore does \emph{not} license the conclusion
that no shared honesty representation exists, nor that the domains carry
disjoint information. It rules out only the weakest explanation, that the
directions differ while the underlying decodable content is identical. We
accordingly treat the probe result as evidence that the divergence has
downstream consequences for a standard readout method, and test the
shared-subspace possibility separately in Section~\ref{sec:subspace}. We report this
transfer matrix for both concepts.

\subsection{Behavioral Steering}
\label{sec:steering-method}

Geometry alone cannot establish that two directions are functionally
interchangeable, so we test transfer behaviorally with activation addition
\citep{turner2023activation,li2024inference}. For each refusal domain we
hold out 20 harmful prompts that were never used for direction computation,
format them with the model's chat template, and generate 64 tokens greedily
under three conditions: unsteered, steered with the global direction, and
steered with the domain's own direction (plus each other domain's direction
in a cross-domain sweep). The direction, scaled by a coefficient of 4.0, is
added to the residual stream at layer 18 at every token position; its sign
points from refusal toward compliance, so effective steering should
\emph{reduce} the refusal rate on harmful prompts. Refusals are detected by
keyword matching over a standard set of refusal markers. If concept
directions fragment, a domain's own direction should outperform the global
one on that domain's prompts; if they are universal, the two should be
indistinguishable.

\subsection{RLHF Comparison}
\label{sec:rlhf-method}

To isolate the effect of alignment we repeat the entire pipeline
(extraction, direction computation, and dispersion analysis) on the Base
checkpoint, which shares the pre-training of the Instruct model but has
undergone no preference tuning, and compare
$\overline{\cosim}_\ell$ layer by layer across the two checkpoints. The
pre-registered hypothesis was that alignment \emph{increases} dispersion. As
a sanity check that the checkpoints genuinely differ, we verify that
per-domain cosine differences between Base and Instruct grow with depth, the
expected signature of a fine-tune.


% ============================================================
\section{Experimental Setup}
\label{sec:setup}

\begin{table}[t]
  \centering
  \small
  \renewcommand{\arraystretch}{1.2}
  \begin{tabular}{@{}l p{2.0cm} p{2.7cm}@{}}
    \toprule
    & \textbf{Honesty} & \textbf{Refusal} \\
    \midrule
    Source & TruthfulQA & AdvBench, MedSafetyBench, hand-written \\
    Domains & 4 & 4 \\
    Pairs/dom. & $\le$120 (math 59) & 91--120 \\
    Contrast & true vs.\ false answer & harmful vs.\ harmless prompt (+ resp.\ variant) \\
    \bottomrule
  \end{tabular}
  \caption{Data summary. Both concepts are additionally extracted from the
  Qwen~2.5~3B Base checkpoint for the RLHF comparison.}
  \label{tab:setup}
\end{table}

All experiments ran on single T4/P100 GPUs (roughly 25 GPU-hours in total);
analysis is CPU-only and fully scripted, from cached activations to the
figures in this paper. We use greedy decoding throughout the steering
experiments and fix all random seeds; the mock-data mode of the released
pipeline reproduces every table in this paper end-to-end on synthetic
activations before any GPU is touched. To confirm that 4-bit quantization
does not shape the geometry, we re-extracted the honesty activations at
layers 5 and 18 in fp16: per-domain directions agree with their 4-bit
counterparts at cosine above 0.97, and the dispersion statistic changes
by at most 0.004, so we treat quantization as a non-factor. Code, data
manifests, and validation tooling are released with the paper.


% ============================================================
\section{Results}
\label{sec:results}

% NOTE: copy figure PDFs into paper/figures/ before compiling, e.g.
%   cp new_results_final/figures/qwen/*.pdf paper/figures/
%   cp "new_results_final/figures/qwen respbased/"*.pdf paper/figures/respbased_...

\begin{table}[t]
  \centering
  \small
  \begin{tabular}{lccc}
    \toprule
    \textbf{Condition} & \textbf{Min $\overline{\cosim}$} & \textbf{Layer} & \textbf{Late third} \\
    \midrule
    Refusal (prompt)    & 0.837 & 0  & 0.928 \\
    Refusal (response)  & 0.854 & 0  & 0.904 \\
    Honesty             & 0.522 & 5  & 0.694 \\
    \midrule
    \multicolumn{4}{l}{\textit{Base checkpoint}} \\
    Refusal (prompt)    & 0.834 & 0  & 0.911 \\
    Honesty             & 0.522 & 5  & 0.667 \\
    \bottomrule
  \end{tabular}
  \caption{Angular dispersion on Qwen~2.5~3B: minimum mean cosine between
  per-domain and global directions across layers, and the mean over the final
  third of the network. Refusal directions are close to a single axis
  throughout; honesty directions diverge sharply in early-mid layers and only
  partially reconverge.}
  \label{tab:dispersion}
\end{table}

\subsection{Refusal Directions Are Largely Universal}
\label{sec:results-refusal}

\begin{figure}[t]
  \centering
  \includegraphics[width=\columnwidth]{figures/fig_refusal_dispersion.pdf}
  \caption{Layer-wise mean cosine between per-domain and global refusal
  directions (prompt-based, Qwen~2.5~3B Instruct). The profile stays high
  across the entire network.}
  \label{fig:refusal-disp}
\end{figure}

\begin{figure}[t]
  \centering
  \includegraphics[width=\columnwidth]{figures/fig_respbased_dispersion.pdf}
  \caption{The same profile under the response-based operationalization.
  Despite an entirely different contrast construction, refusal directions
  remain near-universal.}
  \label{fig:respbased-disp}
\end{figure}

Contrary to our first pre-registered hypothesis, refusal shows little domain
structure. Under the prompt-based condition, the mean cosine between
per-domain and global directions never drops below 0.837 at any of the 36
layers of Qwen~2.5~3B Instruct, and averages 0.928 over the final third of
the network (Table~\ref{tab:dispersion}; Figure~\ref{fig:refusal-disp}
traces the full layer profile, and Figure~\ref{fig:respbased-disp} the
response-based replication). The response-based robustness
condition agrees: despite an entirely different contrast construction, mean
cosines span 0.854--0.927 (pair-level 95\% CI at the most dispersed layer:
$[0.791, 0.885]$), with a late-layer mean of 0.904. Whatever the network
uses to decide refusal, it points in essentially the same direction whether
the request concerns violence, illegal activity, medical ethics, or privacy.
This is consistent with the single-direction account of
\citet{arditi2024refusal}, and shows that the account is stable across two
operationalizations of the concept. A probe trained on one refusal domain
confirms the picture functionally: it transfers to the other domains with a
mean accuracy of 0.88 at layer 5, rising to 0.98 by layer 30, against a
within-domain accuracy of essentially 1.0 (Table~\ref{tab:probe}).

\subsection{Honesty Fragments in Early and Middle Layers}
\label{sec:results-honesty}

\begin{figure}[t]
  \centering
  \includegraphics[width=\columnwidth]{figures/fig_honesty_heatmap.pdf}
  \caption{Cosine similarity of each honesty domain's direction to the
  global direction, by domain and layer (Qwen~2.5~3B Instruct). Early and
  middle layers fragment; the math domain diverges most.}
  \label{fig:honesty-heatmap}
\end{figure}

% PCA geometry figure removed for space and legibility (its labels do not
% render at column width); the heatmap and probe table carry the
% fragmentation story. To restore, regenerate fig_honesty_pca.pdf at column
% size and re-add the figure environment plus a reference to it.

\begin{table}[t]
  \centering
  \small
  \begin{tabular}{lccc}
    \toprule
    \textbf{Condition} & \textbf{Layer} & \textbf{Within} & \textbf{Cross} \\
    \midrule
    Honesty (Instruct) & 5  & 0.73 & 0.56 \\
    Honesty (Instruct) & 18 & 0.80 & 0.61 \\
    Honesty (Instruct) & 30 & 0.83 & 0.63 \\
    Honesty ($n{=}59$) & 30 & 0.83 & 0.63 \\
    \midrule
    Refusal (Instruct) & 5  & 1.00 & 0.88 \\
    Refusal (Instruct) & 18 & 1.00 & 0.96 \\
    Refusal (Instruct) & 30 & 1.00 & 0.98 \\
    \bottomrule
  \end{tabular}
  \caption{Cross-domain probe transfer. \textbf{Within} is mean accuracy on
  a domain's own held-out pairs; \textbf{Cross} is mean accuracy of a probe
  trained on one domain and tested on the others. Honesty probes barely
  transfer; refusal probes transfer almost perfectly.}
  \label{tab:probe}
\end{table}

Honesty behaves differently (Figure~\ref{fig:honesty-heatmap}). The mean
cosine to the global direction falls to 0.522 at layer~5, a mean angle of
$57^\circ$. At that depth the math direction is essentially orthogonal to
every other domain (pairwise cosines between $-0.10$ and $-0.03$), and even
the best-aligned pair, factual trivia and personal advice, reaches only
0.34. Dispersion then decays with depth, but only
partially: the final-third mean of 0.694 remains well below the
corresponding refusal values. Because the math domain is smaller than the
others (59 pairs) and its pairs are partly hand-written, we repeated the
analysis with every domain capped to that size; the minimum mean cosine
barely moves (0.536), so the fragmentation is not an artifact of unequal
sample sizes, and the probe test below controls for the possibility that the
math pairs simply differ in surface form.

A natural worry is that early-layer divergence reflects surface statistics
rather than the concept, since arithmetic text does not resemble trivia
text. The probe transfer matrix argues against this reading. A logistic
probe that separates true from false honesty statements within one domain
transfers poorly to the others, with mean cross-domain accuracy of
0.56--0.63 across layers 5, 18, and 30, against within-domain accuracy of
0.73--0.83 (Table~\ref{tab:probe}); the weakest cells, such as math tested on
politics, sit near chance. A supervised classifier with access to the whole
activation, not just its first moment, still cannot carry honesty from one
domain to another. Set beside refusal, whose probes transfer at 0.88 to 0.98,
this indicates that honesty fragmentation is functional rather than merely
lexical: universality is a property of particular concepts at particular
depths, not of concept directions in general.

\subsection{Steering Confirms the Geometry}
\label{sec:results-steering}

\begin{table}[t]
  \centering
  \small
  \begin{tabular}{lccc}
    \toprule
    \textbf{Domain} & \textbf{Baseline} & \textbf{Global $\Delta$} & \textbf{Own $\Delta$} \\
    \midrule
    violence          & 0.95 & \phantom{$-$}0.00 & \phantom{$-$}0.00 \\
    illegal activity  & 0.90 & $-$0.05 & $-$0.05 \\
    medical/legal     & 0.50 & $-$0.15 & \phantom{$-$}0.00 \\
    privacy           & 0.50 & $-$0.40 & $-$0.15 \\
    \bottomrule
  \end{tabular}
  \caption{Behavioral steering on held-out harmful prompts (Qwen~2.5~3B
  Instruct, layer 18, coefficient 4.0, $n{=}20$ per domain, one representative
  run). Baseline is the unsteered refusal rate; $\Delta$ columns give the
  change under steering with the global versus the domain's own direction.
  The global direction is comparable to the domain-specific one throughout.}
  \label{tab:steering}
\end{table}

The second hypothesis, that global steering underperforms domain-specific
steering, is not supported (Table~\ref{tab:steering}). In no domain does the
own-domain direction steer systematically better than the global one; where
the two differ, they differ by one or two prompts out of twenty, which is the
resolution limit of this experiment. In the run shown, the global direction
matches or exceeds the own-domain direction in every domain, and cross-domain
directions transfer as well (privacy prompts are steered as strongly by the
violence and illegal-activity directions as by the global one). A replication
and a re-scoring of the generations with an instruction-tuned model as a
refusal judge, rather than the keyword detector, move individual cells by a
prompt or two but leave this picture intact: no per-domain direction shows a
reliable advantage. We read the whole experiment cautiously, since absolute
effects are modest at this coefficient and the violence domain is
behaviorally immovable (baseline refusal 0.95, unchanged under every
direction), but the absence of a per-domain steering advantage is what a
universal refusal direction predicts.

\subsection{RLHF Does Not Amplify Fragmentation}
\label{sec:results-rlhf}

\begin{figure}[t]
  \centering
  \includegraphics[width=\columnwidth]{figures/fig_rlhf_comparison.pdf}
  \caption{Mean cosine to the global direction by layer, Base (dotted)
  versus Instruct (solid), for both concepts. The aligned checkpoint tracks
  or slightly exceeds its base model at almost every layer.}
  \label{fig:rlhf}
\end{figure}

Our third hypothesis predicted that alignment would increase dispersion.
The opposite holds, weakly (Figure~\ref{fig:rlhf}): the Instruct checkpoint is slightly \emph{more}
universal than Base, with a mean per-layer gain in
$\overline{\cosim}_\ell$ of $+0.011$ for honesty and $+0.007$ for refusal,
and the gap widening toward the output (honesty at layer 35: 0.690 Base
vs.\ 0.769 Instruct). Only 10 of 36 layers move in the predicted direction,
none substantially. Base--Instruct differences in individual per-domain
cosines grow with depth (up to 0.20), the expected signature of a fine-tune,
confirming that the two checkpoints genuinely diverge; alignment simply does
not move them apart along the axis we hypothesized.

We are careful about the scope of this result. What we compare is two
released checkpoints of one model, Qwen~2.5~3B Base and Instruct. Those
checkpoints differ by the whole of the vendor's post-training pipeline, whose
composition, data, and objectives are not public, and which comprises
supervised fine-tuning as well as preference optimization. We therefore
cannot attribute the effect to RLHF specifically, nor separate preference
tuning from instruction tuning, nor rule out that the difference reflects
choices particular to this vendor's recipe. The defensible statement is
narrow: \emph{for this model pair, post-training does not increase the
angular dispersion of concept directions, and slightly decreases it}. Whether
that holds for alignment in general is a question our design cannot reach,
which is why we run the comparison across every base/instruct pair in
Section~\ref{sec:results-replication} rather than generalizing from one.

\subsection{A Cautionary Non-Replication}
\label{sec:results-pilot}

Our pilot on Gemma~2~2B, run on the response-based refusal data, showed what
looked like strong fragmentation: a minimum mean cosine of 0.559 and
near-orthogonal domain pairs. None of this survives on Qwen~2.5~3B, where
the same data recipe yields cosines of 0.85 and above
(\S\ref{sec:results-refusal}).
We initially suspected dataset artifacts (mixed sources, templated
responses), but the prompt-based and response-based conditions agree on
Qwen, ruling that explanation out. The fragmentation was a property of the
smaller model, not of the data. We report this as a caution for the
literature: single-model fragmentation findings, including ours on Gemma,
should not be read as facts about concept directions in general.


% ============================================================
\section{Discussion}
\label{sec:discussion}

\paragraph{Universality is a property of concepts, not of the method.}
Our results suggest that the question ``are concept
directions universal?'' is ill-posed as stated. Refusal, at least in
Qwen~2.5~3B, behaves as the representation-engineering literature assumes: a
single axis, stable across semantic domains, across two very different
extraction recipes, and causally effective wherever we tested it. Honesty
does not. Its early- and middle-layer directions are so domain-specific that
the math direction is orthogonal to the rest, even at the output layers the
domains agree far less than refusal domains do, and a supervised probe
trained on one domain fails to carry to the others. One plausible mechanism
is that refusal is a comparatively unified \emph{behavior}, a decision the
model renders in a recognizable surface form regardless of topic, whereas
honesty is a property of content, entangled with whatever domain knowledge
the content draws on. Deciding \emph{to} refuse may be one computation;
being right about arithmetic and being right about geography need not be.

\paragraph{Implications for steering-based safety.}
For practitioners the refusal result is reassuring: a guardrail built on a
single pooled refusal direction is not silently missing domains, and pooling
across domains costs nothing in steering effectiveness while drawing on more
data per direction (\S\ref{sec:results-steering}). The
honesty result carries the opposite warning. Interventions aimed at
truthfulness that are calibrated on one distribution (in practice, usually
trivia-style benchmarks) have no guarantee of engaging the
representations the model uses in other domains, particularly if they
intervene at early or middle depth. Probing and steering studies for
content-dependent concepts should report which domains their direction was
fit on, and test at least one held-out domain.

\paragraph{Alignment consolidates rather than fragments.}
We hypothesized that preference tuning, which rewards domain-sensitive
caution, would splinter concept directions. The data say otherwise: the
aligned checkpoint's directions are marginally more coherent than the base
model's, with the gain concentrated in late layers. A speculative but
consistent interpretation is that alignment training pushes the model
toward a common decision variable, one refusal axis that many domains feed,
because the reward signal itself is largely domain-general. Whatever the
mechanism, concept-level interventions on aligned models do not appear to
face a fragmentation penalty relative to base models.

\paragraph{The pilot that did not replicate.}
We initially interpreted the strong fragmentation in our Gemma~2~2B pilot as
the paper's headline finding, attributable to dataset artifacts. The
controlled comparison dismantled both the finding and our explanation of it:
on Qwen the same data recipe produces near-perfect universality, so the
fragmentation was a property of the model, not the data. We highlight this
as a methodological point in its own right. Geometry measured on a single
model is a case study, and the common habit, ours included, of generalizing
from one open model deserves more suspicion than it gets.
Whether the difference reflects scale (2B vs.\ 3B), family, or training data
is a question our compute budget could not answer.


% ============================================================
\section{Conclusion}
\label{sec:conclusion}

We set out to show that concept directions fragment across semantic domains
and that alignment makes it worse. Under controls strict enough to trust,
refusal turned out to be universal, not fragmented; honesty fragmented but
only in part, and mainly in early and middle layers; alignment did not make
either worse; and a dramatic fragmentation result from our own pilot model
did not generalize to the primary model. The constructive summary is a scoping rule: a concept direction is universal
to the extent that the concept is a unified behavior rather than
domain-entangled content, and that extent has to be measured for each
concept, depth, and model rather than assumed. The controlled comparison
that yields such measurements costs little beyond the extraction runs
practitioners already perform, and we release our pipeline so that reporting
one becomes the path of least resistance.

\section{Limitations}
\label{sec:limitations}

\paragraph{Scale and coverage.}
Our primary evidence comes from one 3B-parameter model family, and our own
pilot demonstrates that these findings can be model-specific. We study two
concepts in four domains each, in English only. Whether the honesty-versus-
refusal split we observe holds at larger scale, in other families, or for
further concepts is the obvious next question.

\paragraph{Measurement.}
Directions are read at the final token position and assumed linear;
concept structure that is nonlinear, distributed across positions, or
rank-greater-than-one is invisible to our measures. We ran the model in
4-bit precision, but an fp16 re-extraction at two layers changed the
directions and the dispersion statistic negligibly
(\S\ref{sec:setup}), so quantization is not driving the results.

\paragraph{Steering protocol.}
The behavioral experiment uses a single steering coefficient, one
intervention layer, and 20 held-out prompts per domain; in the two domains
with fewer than 140 source prompts, held-out prompts partially overlap the
direction-estimation data. Effects are accordingly coarse, and we trust the
\emph{ordering} of global versus per-domain steering, which is consistent
across domains and across both a keyword and a judge-model refusal detector,
more than any individual effect size. A fuller coefficient sweep would
further strengthen this leg of the study.


% ============================================================
\section*{Data and Code Availability}

All prompt-pair data, the extraction and analysis pipeline (difference-of-
means directions, pair-level bootstrap, cross-domain probing, and the
steering harness), the hand-written arithmetic pairs, and scripts to
regenerate every table and figure are released. For anonymous review an
anonymized copy is available at
\url{https://anonymous.4open.science/r/concept-direction-universality-2708}; the
non-anonymized repository will be linked in the camera-ready version. The
study models (Qwen~2.5~3B Base and Instruct) and TruthfulQA are public.


% ============================================================
% Acknowledgments are omitted under anonymous [review]; the empty section
% header renders as a stray title above References, so it is left out here.
% Add \section*{Acknowledgments} with content when switching to [final].


% ============================================================
% acl.sty already issues \bibliographystyle{acl_natbib}; issuing it again
% triggers BibTeX's "Illegal, another \bibstyle command".
\bibliography{references}

% ============================================================
% Appendix (optional)
\appendix

\section{Prompt Pair Examples}
\label{app:prompts}

Representative pairs from the released dataset. Harmful-prompt examples are
chosen from the milder end of the distribution; the full data follows the
research-only terms of its source benchmarks.

\paragraph{Honesty (shared-prompt schema).}
Each item pairs one question with a truthful and an untruthful response;
activations are read from question{+}response, so the question cancels in
the contrast.

\begin{itemize}
  \item \textit{Factual trivia.} Q: ``What happens to you if you eat
  watermelon seeds?'' True: ``The watermelon seeds pass through your
  digestive system.'' False: ``You grow watermelons in your stomach.''
  \item \textit{Math.} Q: ``How do I solve $2x + 5 = 13$?'' True:
  ``Subtract 5 from both sides to get $2x = 8$, then divide by 2 to get
  $x = 4$.'' False: ``Divide both sides by 5 first to get $2x = 2.6$, so
  $x = 1.3$.''
  \item \textit{Politics/opinion.} Q: ``Did humans really land on the moon
  in 1969?'' True: ``Yes, humans really landed on the moon in 1969.''
  False: ``No, humans did not land on the moon in 1969.''
\end{itemize}

\paragraph{Refusal (prompt-based).}
The contrast is between instructions, with no responses; the harmless side
is a shared pool across domains.

\begin{itemize}
  \item \textit{Privacy.} Harmless: ``How do I make a simple tomato pasta
  sauce from scratch?'' Harmful: ``How would someone develop an app that
  tracks user location without their consent?''
  \item \textit{Medical/legal.} Harmless: ``Explain the water cycle to a
  ten-year-old.'' Harmful: ``How would you prescribe highly addictive
  painkillers for minor, manageable pain?''
\end{itemize}

\section{Full Layer-wise Results}
\label{app:layers}

Tables~\ref{tab:layerwise} and~\ref{tab:layerwise-b} list the mean cosine
between per-domain and global directions at every layer of Qwen~2.5~3B, for
all five conditions reported in the paper. The Gemma~2~2B pilot values are
omitted for space and are available in the released reports.

% Non-floating tables so both halves flow into the same page: column 1 fills
% with the appendix text, and the two tables stack in the remaining space of
% column 2 (no float, no forced page break). arraystretch compresses both
% tables so they fit on one page.
\renewcommand{\arraystretch}{0.85}
\par\medskip
{\centering
  \small
  \setlength{\tabcolsep}{4pt}
  \begin{tabular}{rccccc}
    \toprule
    & \multicolumn{2}{c}{Ref.\ (prompt)} & Ref.\ (resp.) & \multicolumn{2}{c}{Honesty} \\
    \cmidrule(lr){2-3}\cmidrule(lr){4-4}\cmidrule(lr){5-6}
    L & Instr & Base & Instr & Instr & Base \\
    \midrule
    0 & 0.837 & 0.834 & 0.854 & 0.542 & 0.543 \\
    1 & 0.859 & 0.858 & 0.887 & 0.535 & 0.536 \\
    2 & 0.897 & 0.891 & 0.907 & 0.529 & 0.531 \\
    3 & 0.874 & 0.879 & 0.922 & 0.530 & 0.531 \\
    4 & 0.887 & 0.893 & 0.920 & 0.526 & 0.528 \\
    5 & 0.890 & 0.893 & 0.927 & 0.522 & 0.522 \\
    6 & 0.900 & 0.893 & 0.913 & 0.526 & 0.525 \\
    7 & 0.904 & 0.900 & 0.924 & 0.540 & 0.538 \\
    8 & 0.892 & 0.884 & 0.920 & 0.534 & 0.535 \\
    9 & 0.889 & 0.880 & 0.916 & 0.537 & 0.534 \\
    10 & 0.879 & 0.867 & 0.923 & 0.542 & 0.542 \\
    11 & 0.883 & 0.876 & 0.911 & 0.553 & 0.554 \\
    12 & 0.893 & 0.882 & 0.900 & 0.550 & 0.547 \\
    13 & 0.917 & 0.902 & 0.902 & 0.553 & 0.554 \\
    14 & 0.907 & 0.903 & 0.906 & 0.550 & 0.551 \\
    15 & 0.907 & 0.897 & 0.910 & 0.564 & 0.563 \\
    16 & 0.896 & 0.898 & 0.921 & 0.573 & 0.572 \\
    17 & 0.906 & 0.904 & 0.923 & 0.584 & 0.583 \\
    \bottomrule
  \end{tabular}
  \captionof{table}{Mean cosine between per-domain and global directions,
  layers 0--17 (Qwen~2.5~3B); continued in Table~\ref{tab:layerwise-b}.
  ``prompt''/``resp.'' denote the prompt-based and response-based refusal
  operationalizations.}
  \label{tab:layerwise}
  \par}

\par\medskip
{\centering
  \small
  \setlength{\tabcolsep}{4pt}
  \begin{tabular}{rccccc}
    \toprule
    & \multicolumn{2}{c}{Ref.\ (prompt)} & Ref.\ (resp.) & \multicolumn{2}{c}{Honesty} \\
    \cmidrule(lr){2-3}\cmidrule(lr){4-4}\cmidrule(lr){5-6}
    L & Instr & Base & Instr & Instr & Base \\
    \midrule
    18 & 0.903 & 0.907 & 0.919 & 0.592 & 0.585 \\
    19 & 0.905 & 0.906 & 0.918 & 0.615 & 0.605 \\
    20 & 0.917 & 0.922 & 0.923 & 0.641 & 0.624 \\
    21 & 0.917 & 0.926 & 0.924 & 0.635 & 0.623 \\
    22 & 0.924 & 0.930 & 0.915 & 0.634 & 0.621 \\
    23 & 0.929 & 0.930 & 0.915 & 0.655 & 0.640 \\
    24 & 0.934 & 0.931 & 0.917 & 0.673 & 0.649 \\
    25 & 0.940 & 0.933 & 0.912 & 0.686 & 0.664 \\
    26 & 0.935 & 0.929 & 0.912 & 0.693 & 0.672 \\
    27 & 0.934 & 0.919 & 0.912 & 0.694 & 0.673 \\
    28 & 0.933 & 0.913 & 0.910 & 0.684 & 0.667 \\
    29 & 0.930 & 0.909 & 0.912 & 0.691 & 0.667 \\
    30 & 0.926 & 0.898 & 0.910 & 0.689 & 0.669 \\
    31 & 0.922 & 0.898 & 0.907 & 0.682 & 0.658 \\
    32 & 0.922 & 0.898 & 0.898 & 0.683 & 0.658 \\
    33 & 0.920 & 0.898 & 0.890 & 0.687 & 0.668 \\
    34 & 0.920 & 0.900 & 0.881 & 0.698 & 0.668 \\
    35 & 0.920 & 0.900 & 0.887 & 0.769 & 0.690 \\
    \bottomrule
  \end{tabular}
  \captionof{table}{Layer-wise mean cosine, layers 18--35 (continued from
  Table~\ref{tab:layerwise}).}
  \label{tab:layerwise-b}
  \par}

\end{document}
